% !TEX root = ../macro.tex
\section{Labor Markets and Unemployment}\label{sec:labor-markets}

% BEGIN ORIGINAL LABOR-MARKET OUTLINE (comparison only)
% \plannedsection{Labor Markets and Unemployment}\label{sec:labor-markets}
%
% \subsection{Employment, unemployment, and participation}
% Define the margins hidden by a single aggregate labor-supply variable.
% Distinguish changes in hours from movements into and out of employment.
%
% \subsection{Search, matching, and wage determination}
% Introduce a basic matching framework to explain vacancies and unemployment.
% State how wage setting and hiring incentives determine job creation.
%
% \subsection{Adjustment and persistence}
% Compare search frictions with nominal wage rigidity. Examine which mechanisms
% can make employment recover slowly and connect a recession to longer-run output.
% END ORIGINAL LABOR-MARKET OUTLINE

\subsection{Why can a willing worker and an open job remain apart?}

Imagine a district where a warehouse advertises jobs while nearby residents
remain unemployed. One resident has not seen the advertisement; another needs
a different shift; a third is waiting for an interview. The employer also needs
time to screen applicants and learn whether a match will work. A posted wage
does not instantly place the right person in the right job.

The previous chapters made useful decisions about this problem. RBC described
aggregate labor as a quantity households and firms could coordinate through a
real wage. The wage-contract NK model distinguished the compensation desired
by labor suppliers from the wage employers actually paid. Neither construction
counted the workers waiting to meet an employer. We now make that waiting
process explicit.

Start with people, rather than rename a fall in aggregate hours as unemployment.
Then count movements into and out of jobs. To explain those movements, we need
workers searching, firms recruiting, and a technology that turns their efforts
into matches. Recruiting costs resources before a job produces anything.
The expected value of the eventual match determines whether a firm opens a
vacancy. Once the parties meet, their outside options help determine how they
divide the gain from working together.

This sequence is a teaching reconstruction, not an account of any author's
private thought process. It leads to a basic Diamond--Mortensen--Pissarides
(DMP) search-and-matching model. We use homogeneous jobs and exogenous
separations to isolate job creation and unemployment flows. Endogenous job
destruction, participation choices, capital, and nominal contracts require
additional decisions.\footnote{For a primary-source formulation and its
quantitative limitations, see Shimer (2005),
\href{https://home.uchicago.edu/~shimer/wp/published/uv.pdf}{\textit{The Cyclical
Behavior of Equilibrium Unemployment and Vacancies}}, American Economic Review
95(1), 25--49. The present deterministic examples and parameter choices are
our own; they do not reproduce that paper's calibration or empirical estimates.}

\Needspace{7\baselineskip}
\begin{greenblock}
\textbf{Three different coordination problems.}
A flexible wage coordinates desired hours in the RBC benchmark. An inherited
wage quote restricts nominal adjustment in the previous NK chapter. A matching
friction delays the formation of a working relationship. Unemployment can
therefore persist even when the wage of every newly formed match is negotiable.
\end{greenblock}

\subsection{Count people, hours, and transitions separately}

\subsubsection{An unemployment rate has a denominator}

Let $P$ be the relevant population, $E$ the number employed, $U$ the number
unemployed, and $O$ the number outside the labor force. Define
\begin{align}
 P&=E+U+O, & L&=E+U,\\
 u&=\frac UL, & \ell&=\frac LP, & e&=\frac EP.
\end{align}
The unemployment rate uses the labor force, not the whole population. An
unemployed person is searching and available under the statistical definition
being used; a nonemployed person need not be classified as unemployed.
Exact survey rules must be checked before using data.\footnote{For one explicit
statistical implementation, see the US Bureau of Labor Statistics,
\href{https://www.bls.gov/cps/definitions.htm}{\textit{Concepts and Definitions
(CPS)}}. Its population universe, search criteria, temporary-layoff provisions,
and treatment of absence from work are measurement rules, not assumptions of
the matching model below. No current labor-market estimates are used here.}
Multiplying the two relevant shares gives
\begin{align}
 e&=\frac LP\frac EL
   =\ell\frac{L-U}{L}
   =\ell(1-u).
\end{align}
Thus a lower unemployment rate can accompany a lower employment-population
ratio if participation falls sufficiently.

For a hypothetical population of 100, first take $(E,U,O)=(60,5,35)$ and then
$(58,3,39)$. The two observations imply
\begin{align}
 (\ell_0,u_0,e_0)&=(.65,5/65,.60),\\
 (\ell_1,u_1,e_1)&=(.61,3/61,.58).
\end{align}
Measured unemployment falls from about $7.69\%$ to $4.92\%$, while employment
also falls. These counts alone do not reveal which individuals moved between
states; individual transition data would supply that information.

\subsubsection{Fewer workers and shorter shifts are different adjustments}

Let $h$ denote average physical hours per employed person, measured consistently
with the employment count. Total physical hours satisfy
\begin{align}
 H&=Eh=P\ell(1-u)h,\\
 \log H&=\log P+\log\ell+\log(1-u)+\log h,\\
 d\log H&=d\log P+d\log\ell-\frac{du}{1-u}+d\log h.
\end{align}
The middle identity also holds in exact log differences between two dates.
The last line is its differential, not an exact finite-change formula using
the initial unemployment rate. A firm cutting every shift by one hour changes
the intensive margin. A firm dismissing workers changes the extensive margin.
An aggregate $N$ in a production function does not identify which happened.
Moreover, the CES labor composite in the preceding wage chapter need not equal
physical $H$ when relative wages differ.

\subsubsection{A stock can be stable while many people change state}

For now close participation and normalize $L=1$, so $n=1-u$ is employment.
During a short interval $dt$, an employed worker separates at rate $s>0$ and an
unemployed worker finds a job at rate $f\geq0$. Inflows minus outflows give
\begin{align}
 u(t+dt)-u(t)&=s[1-u(t)]dt-f(t)u(t)dt+o(dt),\\
 \frac{u(t+dt)-u(t)}{dt}&=s[1-u(t)]-f(t)u(t)+\frac{o(dt)}{dt},\\
 \dot u(t)&=s[1-u(t)]-f(t)u(t).
 \label{eq:lm-unemployment-flow}
\end{align}
The dot means a derivative with respect to continuous time. The approximation
keeps one transition in $dt$; the probability of two transitions is of smaller
order. Equal inflows and outflows keep $u$ constant without keeping the same
people unemployed. A recession can raise unemployment through more separations,
slower job finding, or both. A stock series alone cannot separate these routes.

\subsection{The search economy: people, resources, and timing}

The conditions below define a real labor-market benchmark. We temporarily set
aside the NK consumption Euler equation, differentiated products, and Calvo
contracts. A complete NK--matching model would have to specify how these
pieces fit together; its equations cannot be obtained by appending an
unemployment identity to the previous chapter.

\Needspace{12\baselineskip}
\begin{modelsetup}[1]
\begin{modelcondition}{01}{People, jobs, and the employment margin}
There is a unit mass of infinitely lived workers and a continuum of potential
single-worker jobs. Workers are either employed, $n=1-u$, or unemployed and
searching, $u$. Labor-force participation is fixed at one. Each employed worker
supplies one fixed unit of time; hours per employee, search effort, and
on-the-job search are not chosen. Jobs and workers are homogeneous before
matching. There is no population growth, migration, capital, government,
foreign sector, or nominal unit of account.
\end{modelcondition}
\end{modelsetup}

\begin{modelsetup}[1]
\begin{modelcondition}{02}{Continuous time, valuation, and risk}
Time is $t\geq0$, measured in months in the numerical examples. All agents value
flows linearly in consumption units and discount at $r>0$:
$E_t\int_t^\infty e^{-r(\tau-t)}c(\tau)d\tau$.
There is no consumption curvature or additional work disutility. Asset claims
are valued at the same real rate $r$; borrowing constraints and income effects
on search are absent. Individual meetings and separations are random. A law
of large numbers makes aggregate flows deterministic conditional on aggregate
fundamentals. The transition experiments specify an unexpected change at
$t=0$ followed by known deterministic paths, not a stochastic DSGE solution.
\end{modelcondition}
\end{modelsetup}

\begin{modelsetup}[1]
\begin{modelcondition}{03}{Output, home production, and recruitment costs}
An occupied job produces $z>0$ goods per unit time; an unemployed worker
produces $b\geq0$ home goods. All goods are perfect substitutes for consumption.
The baseline has $z>b$. Posting one vacancy costs $\kappa>0$ goods per unit
time while it remains unfilled. There is no one-time posting cost, training
cost, or additional operating fixed cost. If $v$ is vacancies per worker,
market output is $Y=zn$ and aggregate consumption is
\[
 C=zn+bu-\kappa v>0.
\]
Here $b$ is real nonmarket production, not a government-financed unemployment
benefit. A transfer would require a tax counterpart and a different resource
interpretation. Workers own the firms, receive their net profits, and fund
recruitment through those ownership accounts.
\end{modelcondition}
\end{modelsetup}

\begin{modelsetup}[1]
\begin{modelcondition}{04}{Meetings and the two sides of the market}
New matches per unit time satisfy
\[
 m(u,v)=\mathcal A_m u^\xi v^{1-\xi},\qquad
 \mathcal A_m>0,\quad0<\xi<1.
\]
For $u,v>0$, tightness is $\vartheta=v/u$. The worker's job-finding hazard is
$f=m/u$ and a vacancy's filling hazard is $q=m/v$. Matching is random and
undirected; agents take aggregate hazards as given. There is no queue choice
or observable match-quality heterogeneity. The parameter $\xi$ is the
elasticity with respect to searching workers. $\vartheta$ is a vacancy-worker
ratio, unrelated to the Calvo survival probabilities of the preceding chapter.
Hazards are positive rates; they need not be less than one.
\end{modelcondition}
\end{modelsetup}

\begin{modelsetup}[1]
\begin{modelcondition}{05}{Separation, re-entry, and information}
An occupied job disappears at an exogenous Poisson rate $s>0$. The worker
returns to unemployment and the firm loses the match asset. A firm can open
a new vacancy at the prevailing cost. Existing matches have no productive
heterogeneity, so there is no endogenous destruction threshold in the baseline.
Agents observe current fundamentals and share the deterministic forecasts
specified for an experiment. Simultaneous events have zero probability in
continuous time. A newly separated worker can search immediately, but a second
transition within $dt$ has probability of order $dt^2$.
\end{modelcondition}
\end{modelsetup}

\begin{modelsetup}[1]
\begin{modelcondition}{06}{Values, ownership, and finance}
$\mathcal U$ is a worker's unemployment value, $\mathcal W$ the employment
value, $J$ the firm's value of an occupied job, and $V$ the value of a vacancy.
These are discounted lifetime stocks, not flow wages or current output.
The worker values track labor and home income; common portfolio distributions
are omitted from both states. Linear utility and common ownership make that
income additive, so it cancels from the employment surplus and bargaining.
Let $D=\mathcal W-\mathcal U$ and $S=D+J$ when free entry gives $V=0$.
All payments and profits use real goods units. A real flow wage $w$ transfers
income from the firm to its employee; it is not a resource cost for the
economy as a whole. Recruitment $\kappa v$ is a resource cost. Unlimited
potential entry and frictionless funding remove a financial obstacle to
vacancy creation, but they do not remove the matching delay.
\end{modelcondition}
\end{modelsetup}

\begin{modelsetup}[1]
\begin{modelcondition}{07}{How a match divides its gain}
Workers have Nash bargaining weight $0<\eta<1$. A disagreement leaves the
worker unemployed and the firm with its vacancy value, zero under free entry.
At given aggregate outside opportunities, the parties maximize
$D^\eta J^{1-\eta}$ over a transferable positive match surplus. They can
renegotiate the real wage continuously; no wage floor, nominal contract,
private information, or commitment friction is imposed. Both participation
surpluses must be positive. The later fixed-real-wage experiment replaces
this bargaining rule explicitly; it is not the same as a Calvo nominal wage.
\end{modelcondition}
\end{modelsetup}

\begin{modelsetup}[1]
\begin{modelcondition}{08}{Entry, matching, and aggregate accounts}
Vacancies can be opened or withdrawn immediately. In an interior equilibrium
$v>0$, free entry imposes $V=0$ and $\kappa=qJ$. The analysis solves this
positive-entry case; eliminating the gain from a match would require
reconsidering entry and the decision to continue existing jobs.
Matches employ one previously unemployed worker
and fill one vacancy; separations destroy one occupied job. The two sides
therefore use the same aggregate flow $m=fu=qv$. Total firm net cash flow is
$(z-w)n-\kappa v$. Wages plus home output plus that cash flow equal $C$ in
Card 03. Claims between residents have zero aggregate net supply; firm equity
is owned by workers. No independent monetary or aggregate-demand rule is
needed for this risk-neutral real benchmark.
\end{modelcondition}
\end{modelsetup}

\begin{modelsetup}[1]
\begin{modelcondition}{09}{Initial stocks and terminal selection}
Initial unemployment $u(0)\in(0,1)$ is given. Employment cannot jump because
meeting and separation hazards are finite. Vacancies and match asset values
can jump after news. The stationary derivation holds $z,b,s,\mathcal A_m,
\kappa,r,\eta,\xi$ constant. A transition instead solves forward-looking match
values with no explosive terminal component and then evolves $u$ forward.
Discounted terminal assets vanish. The illustrated returning paths converge
to the original stationary equilibrium, and numerical terminal conditions are
checked at successively longer horizons. The maintained interior solution has
$J,D,S>0$ and positive consumption throughout.
\end{modelcondition}
\end{modelsetup}

\begin{modelsetup}[1]
\begin{modelcondition}{10}{Scope of the welfare and extension comparisons}
The planner faces the same matching technology, separation risk, population,
home output, recruitment cost, and initial employment as private agents.
It values total consumption linearly and chooses vacancies; transfers cancel.
The efficiency comparison adds no taxes or additional distortions. Three-state
population accounting later permits participation flows as an identity, but
does not supply a participation optimization problem. Skill loss, mismatch,
endogenous separations, and a nominal NK embedding are identified as further
mechanisms; they are not silently assumed in the baseline solution.
\end{modelcondition}
\end{modelsetup}

\subsection{Matching determines waiting times and unemployment flows}

\subsubsection{The same meeting has two different probabilities}

Substitute $v=\vartheta u$ in the matching technology:
\begin{align}
 m(u,v)&=\mathcal A_m u^\xi(\vartheta u)^{1-\xi}
        =\mathcal A_m u\vartheta^{1-\xi},\\
 f(\vartheta)&=\frac mu=\mathcal A_m\vartheta^{1-\xi},\\
 q(\vartheta)&=\frac mv
   =\frac{\mathcal A_m u\vartheta^{1-\xi}}{\vartheta u}
   =\mathcal A_m\vartheta^{-\xi},\\
 f(\vartheta)&=\vartheta q(\vartheta).
\end{align}
The derivatives make the congestion effects explicit:
\begin{align}
 f'(\vartheta)&=(1-\xi)\mathcal A_m\vartheta^{-\xi}>0,\\
 q'(\vartheta)&=-\xi\mathcal A_m\vartheta^{-\xi-1}<0,\\
 \frac{\vartheta f'}f&=1-\xi,
 &\frac{\vartheta q'}q&=-\xi.
\end{align}
More vacancies per searching worker help that worker find a job, while making
each individual vacancy harder to fill. A hospital competing with many other
recruiters may offer abundant opportunities to an applicant and still wait
months to fill its own post. There is no contradiction between those two
perspectives.

With a constant hazard $f$, let $B(a)$ be the probability that an unemployed
worker has not yet found a job after duration $a$. Then
\begin{align}
 B(a+da)&=B(a)[1-f\,da]+o(da),\\
 B'(a)&=-fB(a),\qquad B(0)=1,\\
 B(a)&=e^{-fa},\qquad \Pr(T\leq a)=1-e^{-fa},\\
 E[T]&=\int_0^\infty B(a)da
     =\left[-\frac{e^{-fa}}f\right]_0^\infty=\frac1f.
\end{align}
The last expression is a mean waiting time to the first job, not the probability
of being employed at the end of a month. That distinction matters because the
worker can find a job and separate again before the month ends.

\Needspace{7\baselineskip}
\begin{blueblock}
\textbf{A rate is not a probability.}
A hazard of $1.2$ per month does not mean a $120\%$ probability. With a constant
hazard, the probability of at least one arrival within a month is
$1-e^{-1.2}$. A final-state transition probability additionally accounts for
later transitions back. Continuous-time rates and discrete-time probabilities
must be converted before comparing their numerical values.
\end{blueblock}

\subsubsection{Solve the stock equation before calling unemployment persistent}

For constant $f$ and $s$, rewrite \eqref{eq:lm-unemployment-flow} as
\begin{align}
 \dot u+(s+f)u&=s, & \bar u&=\frac{s}{s+f},\\
 \dot u&=-(s+f)(u-\bar u).
\end{align}
Let $d(t)=u(t)-\bar u$. Multiplying $\dot d+(s+f)d=0$ by the integrating factor
$e^{(s+f)t}$ gives
\begin{align}
 \frac{d}{dt}[e^{(s+f)t}d(t)]&=0,\\
 e^{(s+f)t}d(t)&=d(0),\\
 u(t)&=\bar u+[u(0)-\bar u]e^{-(s+f)t}.
 \label{eq:lm-stock-solution}
\end{align}
Initial unemployment is inherited. Even if vacancy incentives change at once,
the people waiting for jobs move gradually. The convergence half-life is
$\log2/(s+f)$. A larger finding rate both lowers steady unemployment and speeds
adjustment to it. This model does not imply an arbitrarily long recovery when
its own hazards are high.

Set the initial state first to unemployment, $u(0)=1$, and then to employment,
$u(0)=0$. The exact end-of-interval transition probabilities are
\begin{align}
 P_{U\to E}(\Delta)
 &=1-\left[\frac{s}{s+f}+\frac f{s+f}e^{-(s+f)\Delta}\right]
   =\frac f{s+f}[1-e^{-(s+f)\Delta}],\\
 P_{E\to U}(\Delta)
 &=\frac{s}{s+f}[1-e^{-(s+f)\Delta}].
\end{align}
The remaining two probabilities are their row complements. The first-job
probability $1-e^{-f\Delta}$ and the final employment probability above
coincide only if re-separation is excluded. For small $\Delta$, expanding
$e^{-(s+f)\Delta}=1-(s+f)\Delta+o(\Delta)$ gives the familiar first-order
approximations $P_{U\to E}\simeq f\Delta$ and $P_{E\to U}\simeq s\Delta$.

\subsection{Value a match, divide its surplus, and decide whether to recruit}

\subsubsection{Where the asset equations come from}

Begin with a stationary environment. Track the labor/home-income component
of value; common ownership distributions cancel from the state comparison.
An unemployed worker receives $b\,dt$
now, stays unemployed with probability $1-fdt$, and finds a job with probability
$fdt$. Discounting the continuation gives
\begin{align}
 \mathcal U
 &=b\,dt+e^{-r dt}[(1-fdt)\mathcal U+fdt\mathcal W]+o(dt),\\
 &=b\,dt+(1-rdt)[\mathcal U+fdt(\mathcal W-\mathcal U)]+o(dt),\\
 0&=[b-r\mathcal U+f(\mathcal W-\mathcal U)]dt+o(dt),\\
 r\mathcal U&=b+f(\mathcal W-\mathcal U).
 \label{eq:lm-U-bellman}
\end{align}
The product of two terms of order $dt$ is absorbed in $o(dt)$. Applying the
same reasoning to an employed worker yields
\begin{align}
 \mathcal W
 &=w\,dt+e^{-rdt}[(1-sdt)\mathcal W+sdt\mathcal U]+o(dt),\\
 r\mathcal W&=w+s(\mathcal U-\mathcal W).
 \label{eq:lm-W-bellman}
\end{align}
For an occupied job, current profit is $(z-w)dt$ and separation destroys the
match asset:
\begin{align}
 J&=(z-w)dt+e^{-rdt}(1-sdt)J+o(dt),\\
 (r+s)J&=z-w,\qquad J=\frac{z-w}{r+s}.
 \label{eq:lm-job-value}
\end{align}
Thus $z-w$ is a profit flow, while $J$ capitalizes that flow until separation.
A low separation rate can make a modest flow margin valuable.

For a vacancy, the firm pays while it waits:
\begin{align}
 V&=-\kappa dt+e^{-rdt}[(1-qdt)V+qdt J]+o(dt),\\
 rV&=-\kappa+q(J-V).
\end{align}
If opening a vacancy had positive net value, entry would continue. With
positive vacancies, free entry therefore gives
\begin{align}
 V&=0, & \kappa&=qJ, & J&=\frac\kappa q,\\
 z-w&=(r+s)\frac\kappa q.
 \label{eq:lm-job-creation}
\end{align}
The vacancy has zero value, but the filled job has positive value. The reward
from filling the vacancy pays for expected recruitment costs. Zero vacancy
value does not mean zero wages, profits, or productive match surplus.

\subsubsection{Why the worker's outside option is not simply home output}

Subtract \eqref{eq:lm-U-bellman} from \eqref{eq:lm-W-bellman} and write
$D=\mathcal W-\mathcal U$:
\begin{align}
 rD&=w-b-sD-fD,\\
 (r+s+f)D&=w-b,\qquad D=\frac{w-b}{r+s+f}.
 \label{eq:lm-worker-surplus}
\end{align}
The worker who rejects this job does not necessarily remain unemployed forever.
When $f$ is high, another match is easier to find. That prospect raises the
unemployment value and reduces the extra value of this particular job at a
given wage.

The parties divide total surplus $S=D+J$. Hold the aggregate outside
environment fixed while an individual pair negotiates. An extra unit of
present-value transfer raises $D$ and lowers $J$ by the same amount. The Nash
problem is therefore
\begin{align}
 \max_{0<D<S}\quad &D^\eta(S-D)^{1-\eta},\\
 \max_{0<D<S}\quad &\eta\log D+(1-\eta)\log(S-D),\\
 0&=\frac\eta D-\frac{1-\eta}{S-D},\\
 \eta(S-D)&=(1-\eta)D,\\
 D&=\eta S,\qquad J=(1-\eta)S,\qquad
 D=\frac\eta{1-\eta}J.
 \label{eq:lm-nash-sharing}
\end{align}
The second derivative is $-\eta/D^2-(1-\eta)/(S-D)^2<0$, so the interior
solution is the unique maximum.

To implement this split with a wage, rearrange the worker's employment equation
while holding its outside value $\mathcal U$ fixed:
\begin{align}
 (r+s)\mathcal W&=w+s\mathcal U,\\
 (r+s)D&=w-r\mathcal U,\\
 w-r\mathcal U&=\frac\eta{1-\eta}(z-w),\\
 (1-\eta)w-(1-\eta)r\mathcal U&=\eta z-\eta w,\\
 w&=\eta z+(1-\eta)r\mathcal U.
\end{align}
Now impose the aggregate outside option and free entry:
\begin{align}
 r\mathcal U
 &=b+fD
  =b+\frac\eta{1-\eta}fJ
  =b+\frac\eta{1-\eta}\frac{f\kappa}{q}
  =b+\frac\eta{1-\eta}\kappa\vartheta,\\
 w&=(1-\eta)b+\eta(z+\kappa\vartheta).
 \label{eq:lm-nash-wage}
\end{align}
The wage reflects home output, match productivity, and opportunities elsewhere.
The last term does not reimburse the employee for posting a vacancy. It enters
because tightness improves the employee's outside opportunities.

\Needspace{8\baselineskip}
\begin{redblock}
\textbf{Which outside option is held fixed during bargaining?}
An individual pair takes the market's unemployment value as given. First derive
the Nash surplus shares, then impose the equilibrium outside option. If one
first substitutes a stationary $\mathcal U(w)$ and differentiates as if the
pair changed every wage and meeting opportunity in the economy, one solves
a different bargaining problem. Aggregate feedback belongs in equilibrium,
after the individual negotiation.
\end{redblock}

Finally, add the worker and firm surplus equations:
\begin{align}
 rS&=w-b-(s+f)D+z-w-sJ,\\
 (r+s)S&=z-b-fD,\\
 (r+s+\eta f)S&=z-b,\\
 S&=\frac{z-b}{r+s+\eta f}.
\end{align}
Wages cancel because they divide the match's gain. Search opportunities
remain because employment removes a worker from the pool that can form
other matches. Neither private bargaining nor free entry alone determines
the efficient amount of recruitment.

\subsection{Close the equilibrium and read the Beveridge curve}

\subsubsection{One equation determines market tightness}

Insert \eqref{eq:lm-nash-wage} in \eqref{eq:lm-job-creation}:
\begin{align}
 z-[(1-\eta)b+\eta z+\eta\kappa\vartheta]
   &=(r+s)\frac\kappa{q(\vartheta)},\\
 (1-\eta)(z-b)-\eta\kappa\vartheta
   &=(r+s)\frac\kappa{q(\vartheta)},\\
 \frac{\kappa(r+s)}{\mathcal A_m}\vartheta^\xi
       +\eta\kappa\vartheta&=(1-\eta)(z-b).
 \label{eq:lm-tightness-equilibrium}
\end{align}
The first cost term values recruitment while the eventual job may separate.
The second reflects the wage response to outside opportunities. The right
side is the firm's bargaining share of the flow gain over home production.

Define the left side minus the right side as $F(\vartheta)$. Then
\begin{align}
 F(0)&=-(1-\eta)(z-b)<0,\\
 F'(\vartheta)
   &=\frac{\kappa(r+s)\xi}{\mathcal A_m}\vartheta^{\xi-1}
        +\eta\kappa>0,\\
 \lim_{\vartheta\to\infty}F(\vartheta)&=\infty.
\end{align}
Continuity and strict monotonicity give exactly one positive stationary
tightness. Recover the other quantities in sequence:
\begin{align}
 f&=\mathcal A_m\vartheta^{1-\xi}, &q&=\mathcal A_m\vartheta^{-\xi},\\
 \bar u&=\frac{s}{s+f}, &\bar n&=1-\bar u, &\bar v&=\vartheta\bar u,\\
 w&=(1-\eta)b+\eta(z+\kappa\vartheta),
 &J&=\frac\kappa q, &D&=\frac\eta{1-\eta}J,\\
 r\mathcal U&=b+fD,
 &\mathcal W&=\mathcal U+D, &S&=D+J,\\
 Y&=z\bar n,
 &C&=z\bar n+b\bar u-\kappa\bar v.
\end{align}
This is an equilibrium with people waiting and firms recruiting. The same
matching flow must satisfy both $f\bar u=q\bar v$ and
$f\bar u=s\bar n$. Neither side of the labor market disappears from the
accounts.

\subsubsection{How a change in opportunities affects recruitment}

Implicit differentiation of $F(\vartheta)=0$ gives
\begin{align}
 F_\vartheta\,d\vartheta+F_z\,dz&=0,
 &F_z&=-(1-\eta),\\
 \frac{d\vartheta}{dz}&=\frac{1-\eta}{F_\vartheta}>0,\\
 \frac{d\vartheta}{db}&=-\frac{1-\eta}{F_\vartheta}<0,\\
 \frac{d\vartheta}{d\kappa}
 &=-\frac{(r+s)/q+\eta\vartheta}{F_\vartheta}<0,\\
 \frac{d\vartheta}{ds}&=-\frac{\kappa/q}{F_\vartheta}<0,\\
 \frac{d\vartheta}{d\mathcal A_m}
 &=\frac{\kappa(r+s)\vartheta^\xi}
        {\mathcal A_m^2 F_\vartheta}>0.
\end{align}
Each derivative holds the other primitive parameters fixed and compares
stationary equilibria. A more productive occupied job attracts recruitment;
a better nonmarket alternative reduces the gain from matching. More expensive
recruitment or shorter expected job life discourages entry. Better matching
raises both direct meeting efficiency and equilibrium tightness.

At fixed $s$ and $\mathcal A_m$, trace a tightness change into the two stocks:
\begin{align}
 \frac{d\bar u}{d\vartheta}
 &=-\frac{s f'(\vartheta)}{(s+f)^2}<0,\\
 \frac{d\log\bar v}{d\log\vartheta}
 &=1+\frac{d\log\bar u}{d\log\vartheta}
  =1-\frac{(1-\xi)f}{s+f}
  =\frac{s+\xi f}{s+f}>0.
\end{align}
Thus a productivity improvement produces fewer unemployed workers and more
vacancies in this stationary comparison. If the primitive change also moves
$s$ or $\mathcal A_m$, one must include that direct effect; the last derivative
alone no longer determines the vacancy response.

\subsubsection{A curve of balanced flows does not determine job creation}

Set unemployment inflows equal to outflows and use the matching function:
\begin{align}
 s(1-u)&=\mathcal A_m u^\xi v^{1-\xi},\\
 v^{1-\xi}&=\frac{s(1-u)}{\mathcal A_m u^\xi},\\
 v_{BC}(u)&=\left[\frac{s(1-u)}{\mathcal A_m u^\xi}\right]^{1/(1-\xi)},\\
 \log v_{BC}
 &=\frac1{1-\xi}[\log s+\log(1-u)-\log\mathcal A_m-\xi\log u],\\
 \frac{d\log v_{BC}}{du}
 &=\frac1{1-\xi}\left[-\frac1{1-u}-\frac\xi u\right]<0.
\end{align}
This downward relation is the model's steady-flow Beveridge curve. Job
creation separately determines $\vartheta$, hence the ray $v=\vartheta u$.
Their intersection selects the stationary stocks. A negative relationship
between vacancies and unemployment in a diagram is not itself a wage-setting
or vacancy-entry equation.

A lower $\mathcal A_m$ raises the vacancies required to sustain a given
unemployment rate; a higher $s$ does the same. An outward shift therefore
does not uniquely identify skill mismatch. Worse screening, geographic
barriers, or a change in separations can have related aggregate symptoms.
In a transition, $\dot u$ need not be zero, so the economy need not sit on the
steady-flow curve even when the matching technology has not changed.

\Needspace{7\baselineskip}
\begin{blueblock}
\textbf{Vacancies are an investment decision; unemployment is an inherited stock.}
Free entry determines how much recruitment is worth undertaking. The flow
equation determines how quickly that recruitment reaches waiting workers.
A vacancy jump and a gradual unemployment response are compatible. Treating
the Beveridge curve as a restriction at every date would erase this adjustment.
\end{blueblock}

% Dynamic experiments, wage comparisons, and the planner follow below.
